\documentclass[12pt]{article}
\usepackage[utf8]{inputenc}
\usepackage[T1]{fontenc}
\usepackage[margin=1in]{geometry}
\usepackage{setspace}
\usepackage{xcolor}
\usepackage[colorlinks=true, urlcolor=blue!60!black, linkcolor=blue!60!black]{hyperref}
\setlength{\parindent}{2em}
\setlength{\parskip}{1.25em}
\renewcommand{\baselinestretch}{1.0}
\begin{document}
\begin{center}
{\LARGE \textit{Dames \& Moore v. Regan}, 453 U.S. 654 (1981)\par}
\vspace{0.8em}
{\large Case Brief\par}
\end{center}
\vspace{1em}

\section{Memory Jogger}
During the Iran Hostage Crisis, the Supreme Court upheld President Carter's and President Reagan's executive orders nullifying attachments on frozen Iranian assets under the International Emergency Economic Powers Act (IEEPA) and suspending private U.S. lawsuits in favor of an international claims tribunal based on congressional acquiescence and related statutory tenor.

\section{Detailed Case Facts}
On November 4, 1979, Iranian militants seized the United States Embassy in Tehran and took American diplomatic personnel hostage. Ten days later, on November 14, 1979, President Jimmy Carter declared a national emergency under the International Emergency Economic Powers Act (IEEPA), 50 U.S.C. \S\S\ 1701--1706, and issued Executive Order No. 12170, blocking the removal or transfer of all property and interests in property of the Government of Iran, its instrumentalities, and the Central Bank of Iran subject to United States jurisdiction (\href{https://www.courtlistener.com/opinion/110564/dames-moore-v-regan/}{CourtListener}; \href{https://supreme.justia.com/cases/federal/us/453/654/}{Justia}). On November 15, 1979, the Treasury Department's Office of Foreign Assets Control promulgated 31 C.F.R. \S\ 535.203(e), declaring that unless licensed or authorized, any attachment, judgment, lien, execution, or other judicial process against blocked Iranian property was null and void, and 31 C.F.R. \S\ 535.805, stipulating that any license or authorization could be ``amended, modified, or revoked at any time.'' On November 26 and December 19, 1979, the Treasury issued general licenses permitting judicial proceedings and prejudgment attachments against Iranian property, while prohibiting the entry of final judgments or decrees.

On December 19, 1979, petitioner Dames \& Moore, a Los Angeles-based engineering firm, sued the Government of Iran, the Atomic Energy Organization of Iran, and several Iranian banks in the United States District Court for the Central District of California. Dames \& Moore alleged that its wholly owned subsidiary had contracted to conduct site studies for a proposed nuclear power plant in Iran, that the Atomic Energy Organization terminated the contract for convenience on June 30, 1979, and that Iran owed \$3,436,694.30 plus interest for completed services. Relying on the Treasury's general license, the District Court issued orders of prejudgment attachment against the property of certain Iranian banks to secure a potential judgment (\href{https://supreme.justia.com/cases/federal/us/453/654/}{Justia}).

On January 19, 1981, the United States and Iran entered into an executive agreement mediated by Algeria---embodied in the two Algiers Declarations---to secure the release of the 52 American hostages, who were freed the next day, January 20, 1981. Under the Algiers Accords, the United States obligated itself to terminate all legal proceedings in U.S. courts involving claims of U.S. nationals against Iran, nullify all attachments and judgments obtained in those suits, transfer all blocked Iranian assets held in U.S. banks by July 19, 1981, and establish a binding Iran--United States Claims Tribunal backed by a \$1 billion security account deposited in the Bank of England. President Carter immediately issued Executive Orders Nos. 12276--12285 revoking the prior attachment licenses, nullifying post-November 14, 1979 non-Iranian interests in frozen assets, and ordering banks to transfer Iranian funds to the Federal Reserve Bank of New York. On February 24, 1981, President Ronald Reagan issued Executive Order No. 12294 ratifying Carter's orders and ``suspending'' all claims eligible for presentation to the Claims Tribunal, providing that any claim over which the Tribunal lacked jurisdiction would revive in U.S. courts (\href{https://tile.loc.gov/storage-services/service/ll/usrep/usrep453/usrep453654/usrep453654.pdf}{Library of Congress, 453 U.S. 654}).

The case implicated several core statutory and constitutional provisions: Article II executive power and Article III judicial power; the Fifth Amendment Just Compensation Clause; IEEPA \S\ 203, 50 U.S.C. \S\ 1702(a)(1)(B) (derived from \S\ 5(b) of the Trading With the Enemy Act); the Hostage Act of 1868, 22 U.S.C. \S\ 1732; the International Claims Settlement Act of 1949, 22 U.S.C. \S\ 1621 \textit{et seq.}; the Foreign Sovereign Immunities Act of 1976 (FSIA), 28 U.S.C. \S\S\ 1330, 1602 \textit{et seq.}; and the Tucker Act, 28 U.S.C. \S\ 1491.

\section{Procedural History}
On January 27, 1981, Dames \& Moore moved for summary judgment in its underlying contract action against the Government of Iran and the Atomic Energy Organization (but not the Iranian banks). The District Court granted summary judgment for the contract amount plus interest, and Dames \& Moore sought writs of garnishment and execution in Washington state court. However, on May 28, 1981 (amended June 8, 1981), the District Court stayed execution of the judgment pending appeal, vacated all prejudgment attachments against the Iranian defendants, and stayed further proceedings against the bank defendants in light of the Executive Orders (\href{https://supreme.justia.com/cases/federal/us/453/654/}{Justia}).

Meanwhile, on April 28, 1981, Dames \& Moore filed a separate suit in the U.S. District Court for the Central District of California against the United States and Secretary of the Treasury Donald Regan seeking declaratory and injunctive relief to block enforcement of the Executive Orders and Treasury regulations as beyond the President's statutory and constitutional authority. On May 28, 1981, the District Court denied a preliminary injunction and dismissed the complaint for failure to state a claim. Dames \& Moore appealed to the Ninth Circuit on June 3, 1981, and the District Court entered an injunction pending appeal barring transfer of assets subject to petitioner's writs. Because lower courts had reached conflicting conclusions and the Algiers Accords imposed a July 19, 1981 transfer deadline, Dames \& Moore petitioned for certiorari before judgment under 28 U.S.C. \S\ 2101(e). The Supreme Court granted certiorari on an expedited schedule, heard oral argument on June 24, 1981, and issued its decision eight days later on July 2, 1981, affirming the District Court's judgment and ordering the mandate to issue forthwith (\href{https://www.courtlistener.com/opinion/110564/dames-moore-v-regan/}{CourtListener}).

\section{Judicial Votes}
\begin{itemize}
  \item \textbf{Majority (8--1 on Part V; 8--1 on footnote 6; 9--0 on core executive-power judgment):} Justice Rehnquist delivered the opinion of the Court, joined in full by Chief Justice Burger and Justices Brennan, Stewart, White, Marshall, and Blackmun. Justice Stevens joined all parts except Part V. Justice Powell joined all parts except footnote 6.
  \item \textbf{Concurrence in Part:} Justice Stevens filed an opinion concurring in part (joining Parts I--IV, declining to join Part V on Tucker Act jurisdiction).
  \item \textbf{Concurrence in Part and Dissent in Part:} Justice Powell filed an opinion concurring in part and dissenting in part (joining the Court's opinion except footnote 6's holding that nullification of the attachments was not a Fifth Amendment taking).
\end{itemize}

\section{Holding}
\textbf{Judgment:} Affirmed. The District Court's dismissal of Dames \& Moore's complaint was upheld, validating the nullification of attachments, the transfer of blocked Iranian assets, and the suspension of private claims in U.S. courts.

\textbf{Rule Announced:} First, IEEPA \S\ 1702(a)(1)(B) expressly authorizes the President during a declared national emergency to nullify post-freeze judicial attachments and compel the transfer of frozen foreign assets, placing such action in Category One of Justice Jackson's \textit{Youngstown Sheet \& Tube Co. v. Sawyer}, 343 U.S. 579 (1952), framework. Second, although neither IEEPA nor the Hostage Act specifically authorizes suspending private lawsuits in U.S. courts, where the settlement of claims is a necessary incident to resolving a major foreign policy dispute, an alternative adjudicative forum is provided, and Congress has historically acquiesced in executive claims settlement (as evidenced by the tenor of IEEPA, the Hostage Act, and the International Claims Settlement Act), the President possesses implied authority to suspend pending claims without divesting Article III jurisdiction. Third, while nullifying revocable post-freeze attachments does not constitute a Fifth Amendment taking, the possibility that suspending the underlying claims could effect a taking makes ripe the question of Tucker Act jurisdiction, and the treaty exception (28 U.S.C. \S\ 1502) does not bar a takings suit in the Court of Claims (\href{https://supreme.justia.com/cases/federal/us/453/654/}{Justia}).

\textbf{Holding vs. Dicta:} The Court's strict holdings are confined to (1) statutory authorization under IEEPA to revoke post-freeze attachment licenses and transfer blocked assets, (2) implied executive authority to suspend pending claims against a foreign sovereign in favor of an international tribunal where Congress has acquiesced, and (3) the availability of a Tucker Act remedy in the Court of Claims. The Court explicitly emphasized as a limiting caution---not a broad rule---that it was not recognizing ``plenary power'' in the President to settle all claims against foreign entities, and its remarks on the merits of whether claim suspension actually constitutes a taking are reserved dicta because the takings claim itself was unripe.

\section{Analysis of Opinions}

\subsection{Majority Opinion (Justice Rehnquist)}
Justice Rehnquist opened by adopting Justice Jackson's celebrated tripartite concurrence in \textit{Youngstown Sheet \& Tube Co. v. Sawyer}, 343 U.S. 579, 635--638 (1952), while refining it from three rigid ``pigeonholes'' into a ``spectrum running from explicit congressional authorization to explicit congressional prohibition'' (\href{https://supreme.justia.com/cases/federal/us/453/654/}{Justia}). Emphasizing judicial restraint under \textit{Ashwander v. TVA}, 297 U.S. 288 (1936), Rehnquist resolved the case in three distinct steps.

First, addressing the nullification of attachments and transfer of assets (Part III), Rehnquist held that the plain text of IEEPA \S\ 1702(a)(1)(B)---empowering the President to ``direct and compel, nullify, void, prevent or prohibit'' any ``transfer'' or ``exercising any right, power, or privilege'' regarding foreign property---directly authorized President Carter's orders. Rejecting petitioner's argument that IEEPA's 1977 repeal of peacetime ``vesting'' authority under the Trading With the Enemy Act limited the President to merely maintaining a freeze, Rehnquist stressed that frozen assets serve as a diplomatic ``bargaining chip'' that private litigants could not dismantle through revocable, post-freeze attachments. Because the President acted with express congressional authorization, his action enjoyed the ``strongest of presumptions'' under \textit{Youngstown} Category One. In footnote 6, the Court added that because petitioner's attachments were obtained under expressly revocable licenses, petitioner acquired no compensable property interest under the Fifth Amendment.

Second, addressing the suspension of pending claims (Part IV), Rehnquist candidly acknowledged that neither IEEPA (which governs property, not \textit{in personam} lawsuits) nor the 1868 Hostage Act (enacted to protect naturalized citizens from forced foreign repatriation) specifically authorized suspending lawsuits. Nevertheless, relying on the ``general tenor'' of those statutes and an unbroken history of executive claims settlement dating back to the 1799 \textit{Wilmington Packet} case---including at least ten executive settlements after 1952---Rehnquist concluded that Congress had implicitly approved the practice through the International Claims Settlement Act of 1949 and its repeated amendments (\href{https://www.everycrsreport.com/reports/LSB11281.html}{CRS Report LSB11281}). Invoking Justice Frankfurter's \textit{Youngstown} ``gloss'' and \textit{United States v. Midwest Oil Co.}, 236 U.S. 459 (1915), Rehnquist held that long-continued executive practice acquiesced in by Congress raised a presumption of congressional consent. He further rejected petitioner's FSIA and Article III objections, reasoning that suspending claims did not strip federal court jurisdiction but merely changed the substantive rule of law under \textit{United States v. Schooner Peggy}, 1 Cranch 103 (1801), and noting that the Iran--U.S. Claims Tribunal provided a meaningful alternative forum.

Third, in Part V, Rehnquist agreed with both parties that whether the claim suspension effected a Fifth Amendment taking was unripe, but held that the availability of a Tucker Act remedy in the Court of Claims was ripe and unobstructed by the treaty exception in 28 U.S.C. \S\ 1502.

\subsection{Concurrence in Part (Justice Stevens)}
Justice Stevens joined Parts I through IV of the Court's opinion, agreeing fully with the statutory and constitutional validation of the Executive Orders. He declined to join Part V, writing in a two-sentence concurrence that the likelihood that requiring Dames \& Moore to pursue its claim before an alternative international forum would constitute an unconstitutional Fifth Amendment ``taking'' was ``so remote that I would not address the jurisdictional question considered in Part V'' (\href{https://supreme.justia.com/cases/federal/us/453/654/}{Justia}).

\subsection{Concurrence in Part and Dissent in Part (Justice Powell)}
Justice Powell joined the Court's opinion except footnote 6, dissenting from the majority's summary conclusion that nullifying petitioner's attachments could not constitute a compensable Fifth Amendment taking. Powell argued that under \textit{Armstrong v. United States}, 364 U.S. 40 (1960), an attachment entitling a creditor to resort to specific property is a recognized property right, and it was both premature and potentially erroneous to foreclose takings claims across hundreds of pending Iran cases whose facts were not before the Court. Powell emphasized that while the political branches possessed extraordinary foreign-affairs powers to use private claims as ``bargaining chips,'' those powers could not displace the Just Compensation Clause, which prevents the Government from forcing a few private citizens to bear burdens that in fairness belong to the public as a whole (\href{https://supreme.justia.com/cases/federal/us/453/654/}{Justia}).

\section{Examples: Future Applications}
\begin{enumerate}
  \item \textbf{Same-Side Application (Executive Claims Settlement Preempting State Law):} A President negotiates a sole executive agreement with a European government to establish an international compensation foundation for historical expropriation claims by U.S. nationals, and a state enacts a statute compelling foreign insurers to disclose policies in aid of state-court litigation. Under \textit{Dames \& Moore}, the President's long-acquiesced-in authority to resolve foreign claims by executive agreement preempts the conflicting state law, as the Supreme Court later held in \href{https://supreme.justia.com/cases/federal/us/539/396/}{\textit{American Insurance Ass'n v. Garamendi}}, 539 U.S. 396 (2003).
  \item \textbf{Same-Side Application (Revoking Licensed Asset Attachments During Sanctions):} After declaring a national emergency under IEEPA against a hostile foreign regime and freezing its U.S. bank deposits, the Treasury issues a revocable license allowing U.S. contractors to attach those accounts; two years later, to seal a diplomatic peace accord, the President revokes the license and transfers the funds to a reconstruction escrow. Because IEEPA \S\ 1702(a)(1)(B) expressly authorizes nullifying post-freeze interests and compelling transfers, the executive action falls squarely in \textit{Youngstown} Category One and is valid.
  \item \textbf{Opposite-Side Application (Unilateral Presidential Directive to State Courts Under Non-Self-Executing Treaty):} Following an International Court of Justice judgment ordering reconsideration of foreign nationals' state criminal convictions, the President issues a memorandum directing state courts to disregard state procedural default rules without congressional authorization. Unlike \textit{Dames \& Moore}, there is no ``systematic, unbroken, executive practice'' of presidents unilaterally overriding state criminal judgments to enforce non-self-executing treaties; thus, the directive exceeds Article II authority, as held in \href{https://supreme.justia.com/cases/federal/us/552/491/}{\textit{Medell\'in v. Texas}}, 552 U.S. 491 (2008).
  \item \textbf{Opposite-Side Application (Executive Action Defying an Express Statutory Prohibition):} Congress enacts a clear statute prohibiting the transfer of frozen assets of a designated state sponsor of terrorism unless Congress passes a joint resolution of approval, and the President nonetheless orders the funds transferred by executive order to resolve a diplomatic standoff. Because Congress has explicitly prohibited the action rather than acquiesced, the President stands at \textit{Youngstown} Category Three (``lowest ebb''), where \textit{Dames \& Moore}'s congressional-acquiescence rationale cannot rescue the order (cf. \href{https://constitution.congress.gov/browse/essay/artII-S1-C1-8/ALDE_00013797/}{Constitution Annotated on \textit{Zivotofsky v. Kerry}}).
  \item \textbf{Fence-Sitter Application (Broad Emergency Tariffs or Domestic Economic Restructuring Under IEEPA):} A President declares a national emergency over persistent trade deficits or cross-border fentanyl trafficking and invokes IEEPA \S\ 1702(a)(1)(B) to impose sweeping, indefinite global tariffs on imported goods. Supporters argue that \textit{Dames \& Moore} reads IEEPA's verbs (``regulate . . . importation'') expansively and treats related trade statutes as inviting broad executive discretion; opponents argue that tariffs are a core Article I taxing power governed by detailed statutory trade regimes (like Section 232 and Section 301), lacking the narrow historical gloss of claims settlement and triggering major-questions scrutiny (\href{https://www.lawfaremedia.org/article/the-case-against-ieepa-tariffs}{Lawfare}).
\end{enumerate}

\section{Critique}
\textit{Dames \& Moore} is widely celebrated for cementing Justice Jackson's \textit{Youngstown} framework as constitutional orthodoxy, yet scholars across the ideological spectrum have sharply criticized its reasoning as an ``end-run'' around separation of powers driven by the diplomatic exigencies of an eight-day drafting deadline.

\textbf{Logical Weaknesses and Doctrinal Tensions:} The opinion's central logical vulnerability lies in Part IV's alchemy of turning statutory silence into implicit authorization. Having explicitly found that neither IEEPA nor the Hostage Act authorized suspending pending lawsuits, the Court nonetheless relied on their ``general tenor'' to ``invite'' independent presidential action (\href{https://www.everycrsreport.com/reports/LSB11281.html}{CRS Report LSB11281}). Critics point out that IEEPA was enacted in 1977 specifically to \textit{curtail} peacetime presidential emergency powers after the abuses of the Trading With the Enemy Act; treating a power-limiting statute as an invitation to exercise unenumerated powers inverts legislative intent. Similarly, the Court's reliance on ``congressional acquiescence'' assumes that legislative inaction signals consent, ignoring bicameralism, veto-gate dynamics, collective-action barriers, and the practical impossibility of Congress overriding a President during an acute hostage crisis (\href{https://scholarship.law.nd.edu/cgi/viewcontent.cgi?article=4590&context=ndlr}{Notre Dame Law Review}; \href{https://www.gwlr.org/wp-content/uploads/2016/06/84-Geo.-Was.-L.-Rev.-668.pdf}{George Washington Law Review}).

\textbf{Progressive and Institutionalist Critique:} Progressive and civil-liberties scholars argue that \textit{Dames \& Moore} dangerously aggrandized the ``imperial presidency'' in foreign affairs. By transforming \textit{Youngstown}'s Category Two ``zone of twilight'' into a de facto presumption in favor of the Executive whenever Congress has legislated in adjacent fields, the decision weakens legislative checks on war, sanctions, and diplomacy. Critics also echo Justice Powell's concern that extinguishing private judicial remedies and liens without guaranteed full recovery sacrifices individual due process and property rights on the altar of executive expediency, while the Court's insistence that it was merely changing ``substantive law'' rather than stripping Article III jurisdiction relies on a formalistic distinction that masks executive interference with pending federal adjudication.

\textbf{Originalist, Textualist, and Formalist Critique:} Originalist and textualist scholars fault \textit{Dames \& Moore} for subordinating constitutional structure and statutory text to functionalist pragmatism. Under the Treaty Clause (Art. II, \S\ 2, cl. 2) and Bicameralism and Presentment (Art. I, \S\ 7), federal law that overrides existing domestic rights and pending federal causes of action should require Senate ratification or congressional enactment---not a unilateral sole executive agreement (\href{https://arizonalawreview.org/pdf/51-4/51arizlrev1035.pdf}{Arizona Law Review}). Formalists also criticize the Court's reliance on ``historical gloss'' and post-1952 executive practice to construct implied presidential power, warning that allowing executive assertions plus legislative silence to alter the separation of powers creates a one-way ratchet favoring Article II. Indeed, in \href{https://supreme.justia.com/cases/federal/us/552/491/}{\textit{Medell\'in v. Texas}} (2008), a more formalist Court cabined \textit{Dames \& Moore} to its narrow ``claims-settlement'' facts and warned that past executive practice does not by itself create constitutional power.

\section{Key Quotations}
\begin{quote}
``Justice Jackson himself recognized that his three categories represented `a somewhat over-simplified grouping,' 343 U.S. at 635, and it is doubtless the case that executive action in any particular instance falls not neatly in one of three pigeonholes, but rather at some point along a spectrum running from explicit congressional authorization to explicit congressional prohibition.'' --- Justice Rehnquist, Opinion of the Court, 453 U.S. at 669 (\href{https://supreme.justia.com/cases/federal/us/453/654/}{Justia}).
\end{quote}

\begin{quote}
``Such failure of Congress specifically to delegate authority does not, `especially . . . in the areas of foreign policy and national security,' imply `congressional disapproval' of action taken by the Executive.'' --- Justice Rehnquist, Opinion of the Court, 453 U.S. at 678 (\href{https://supreme.justia.com/cases/federal/us/453/654/}{Justia}).
\end{quote}

\begin{quote}
``Past practice does not, by itself, create power, but `long-continued practice, known to and acquiesced in by Congress, would raise a presumption that the [action] had been [taken] in pursuance of its consent. . . .'\,'' --- Justice Rehnquist, Opinion of the Court, 453 U.S. at 686 (\href{https://supreme.justia.com/cases/federal/us/453/654/}{Justia}).
\end{quote}

\begin{quote}
``Finally, we reemphasize the narrowness of our decision. We do not decide that the President possesses plenary power to settle claims, even as against foreign governmental entities.'' --- Justice Rehnquist, Opinion of the Court, 453 U.S. at 688 (\href{https://supreme.justia.com/cases/federal/us/453/654/}{Justia}).
\end{quote}

\begin{quote}
``The Government must pay just compensation when it furthers the Nation's foreign policy goals by using as `bargaining chips' claims lawfully held by a relatively few persons and subject to the jurisdiction of our courts. The extraordinary powers of the President and Congress upon which our decision rests cannot, in the circumstances of this case, displace the Just Compensation Clause of the Constitution.'' --- Justice Powell, concurring in part and dissenting in part, 453 U.S. at 691 (\href{https://supreme.justia.com/cases/federal/us/453/654/}{Justia}).
\end{quote}

\end{document}
